\section{Completion, operations, and the overlap}
\label{complete:section}

We now describe how a twisted periodic complex is recovered from its
restriction to the height-one open and its completion along the
height-two stratum. The fracture square supplies this reconstruction.
To use it for the construction of our object, we also need an explicit
description of finite constructions from the completed chart sheaf.
We obtain one by representing every operation by compatible ordinary
maps on all nilpotent thickenings. Matrices of these maps then form
actual complexes, whose sections and localized sections can be computed.
The last part identifies an intrinsic topology on the overlap when the
original object has finite sections homotopy in every degree, and proves
continuity for every map between such completed objects.

We work at a prime $p>5$, as in the geometric construction. Throughout
this section the categories are the characteristic-$p$ categories
$\mathbf D(\mathcal M_0)$, $\mathbf D(\mathcal M_k)$, and
$\mathbf D(\mathcal N)$ of Section~\ref{geom:periodic-foundations}.
The line section $t=v_1$ has degree $D=2(p-1)$. Recall that the
homological and weight conventions are
$C_{j+D}=C_j(1)$, $d:C_j\to C_{j-1}$, and $C[b]_j=C_{j-b}$,
with the usual sign on the differential of a shift. We write
$\pi_iC=[\mathcal O[i],C]$ for
sections homotopy; this is distinct from the sheaf homology of $C$.
We retain $R=k[[x]]$, $R_d=k[x]/(x^d)$, $\Lambda=\mu_q(k)$, and
$\Delta=R[[H]]$ from the nilpotent charts. All limits and colimits of
objects in the derived categories are homotopy limits and colimits.

\subsection{The fracture square}

For $d\geq 1$ define
\begin{equation}
\label{complete:definitions}
 P/t^d=\cofib\bigl(P[dD]\xrightarrow{t^d}P\bigr),
 \qquad
 \widehat P=\holim_d P/t^d,
 \qquad
 LP=P[1/t].
\end{equation}
The multiplication maps, their cones, and the transition maps in these
towers are chosen naturally at the level of twisted complexes. We call
an object \emph{$t$-invertible} if multiplication
$t:P[D]\to P$ is an equivalence, and \emph{complete} if
$P\to\widehat P$ is an equivalence.

For the general fracture description relating completion and
localization, compare Barthel--Heard--Valenzuela
\cite[Corollaries~2.26 and~2.28]{BHV2018}. We prove the required form
directly for arbitrary twisted periodic complexes in these categories.

\begin{proposition}[Fracture]
\label{complete:fracture}
The functor $L$ is the flat affine restriction and pushforward for the
open $t\ne 0$, whose derived category it embeds fully faithfully.
Completion is exact and idempotent. Its values are right orthogonal,
in every shift, to $t$-invertible objects. For every $P$, the fiber of
$P\to\widehat P$ is $t$-invertible, and there are natural equivalences
\[
 P/t^d\simeq\widehat P/t^d \qquad(d\geq 1).
\]
The following square is a homotopy pullback:
\begin{equation}
\label{complete:fracture-square}
\begin{tikzcd}[column sep=large,row sep=large]
 P \arrow[r] \arrow[d] & \widehat P \arrow[d]\\
 LP \arrow[r] & L\widehat P.
\end{tikzcd}
\end{equation}
Thus the original object is recovered by matching its open and complete
parts over the localized completion. Moreover,
\begin{equation}
\label{complete:localization-sections}
 \pi_i LP=\mathop{\mathrm{colim}}_{s\geq 0}\pi_{i+sD}P.
\end{equation}
\end{proposition}

\begin{proof}
The localization telescope is
$P\to P[-D]\to P[-2D]\to\cdots$, with transition maps induced by $t$.
Exactness of filtered colimits represents this telescope by the ordinary
termwise colimit. On the affine charts this is localization at the
principal open defined by $t$, so it is the asserted open push-pull
functor. Flat affine adjunction identifies the open category with its
essential image.

Taking the inverse limit of the multiplication-cofiber triangles shows
that the fiber of $P\to\widehat P$ is the limit of the tower
\[
 \cdots\xrightarrow{t}P[(d+1)D]
 \xrightarrow{t}P[dD]\xrightarrow{t}\cdots.
\]
Shifting this tower by one position does not change its limit.
Consequently $t$ is invertible on this fiber. On the other hand,
$t^d$ acts nullhomotopically on $P/t^d$, by the standard homotopy on
its multiplication cone. If $A$ is $t$-invertible, naturality of $t$
therefore gives $[A[j],P/t^d]=0$ for all $j$: multiplication by $t^d$
is simultaneously invertible from the source and zero on the target.
The same assertion for mapping spectra, followed by the inverse limit,
gives $[A[j],\widehat P]=0$.

Inverse limits are exact in the stable category, so completion is exact.
If $Q$ is right orthogonal to the $t$-invertible objects, then the fiber
of $Q\to\widehat Q$ is both $t$-invertible and right orthogonal to
the $t$-invertible objects. It is therefore zero. This proves
idempotence and characterizes the complete objects by that
orthogonality. The $t$-invertible fiber of $P\to\widehat P$ has zero
quotient by every $t^d$, giving the asserted equivalence of reductions.
Localizing this same fiber triangle leaves its fiber unchanged, which
proves the pullback assertion in \eqref{complete:fracture-square}.

Finally, the finite cohomological bounds of
Proposition~\ref{geom:cohomological-bounds}, together with exact filtered
colimits, allow sections-homotopy to commute with this telescope.
Equivalently, compute derived sections with the bounded cohomological
range established there and commute the filtered colimit through the
resulting calculation. This gives
\eqref{complete:localization-sections}.
\end{proof}

The fracture square reduces reconstruction to a map on the overlap.
The following description of the open categories specifies which maps
can be chosen there and records a vanishing used in the transfer step.

\begin{lemma}[Objects on the open]
\label{complete:open-splitting}
The category $\mathbf D(\mathcal N[1/t])$ is the category of
$D$-periodic complexes of $k$-vector spaces. Every object is the sum,
over the $D$ residue classes, of its pure homology objects, and every
graded periodic map on homology is realizable.

The category $\mathbf D(\mathcal M_0[1/t])$ is the category of
$D$-periodic complexes of smooth $Z$-representations over $\F_p$.
Every object also splits as the sum of its pure homology objects.
Maps can have an $\Ext^1$ component, but a pure object has no
self-map to its $2$-shift when $D>3$.
\end{lemma}

\begin{proof}
The chart identifications give
$\mathcal N[1/t]=\Spec k$ and
$\mathcal M_0[1/t]=BZ$, and the nowhere-vanishing section $t$
trivializes the periodicity line. The assertion for vector spaces
follows by splitting cycles and boundaries.

Fix the chosen topological generator $c$ of $Z$ and write $s=c-1$.
Smooth representations identify with locally $s$-nilpotent
$\F_p[s]$-modules: the identity
$(1+s)^{p^a}=1+s^{p^a}$ identifies annihilation by a power of $s$
with invariance under a sufficiently small open subgroup. This
torsion category is hereditary. Indeed, embed a torsion module in a
divisible $\F_p[s]$-module and take its $s$-power-torsion submodule.
That submodule is still divisible: division by $s$ can be performed
within the torsion submodule, and every polynomial prime to $s$ acts
invertibly on it. It is thus injective. Its quotient by the original
module is again divisible and torsion, giving an injective resolution
of length one in the torsion category.

Represent a periodic complex by a fibrant complex $I$ with injective
terms. Write $Z_j$, $B_j$, and $H_j$ for its cycles, boundaries, and
homology. From $0\to Z_{j+1}\to I_{j+1}\to B_j\to 0$ and
heredity we obtain
\[
 \Ext^1(H_j,B_j)=\Ext^2(H_j,Z_{j+1})=0.
\]
The identity of $H_j$ therefore lifts to a map $H_j\to Z_j$.
Choose these lifts for one representative of each period residue and
extend them periodically. They give a quasi-isomorphism from the
finite sum of shifted pure $H_j$ to $I$. Finally, periodizing a
length-one injective resolution of a pure object computes its maps;
when $D>3$, neither its resolution degree $0$ nor its resolution
degree $1$ can contribute to a self-map with shift $2$.
\end{proof}

\subsection{The complete operation algebra}

Recall from Proposition~\ref{geom:coinduction} that
$T'=f'_*\mathcal O$ is $t$-monic and that
\[
 T'_d=T'/t^d=f'_*\mathcal O_{\mathcal X_k/(t^d)}
\]
is a pure sheaf. On the nilpotent chart it is
$C_{\mathrm{lc}}(H,R_d)$, coinduced from $\Lambda$, with diagonal
semilinear action $(gF)(h)=g(F(g^{-1}h))$ in coset coordinates.
The chart basis $u$ of the weight line satisfies $t=xu$ and
$h(u)=x/h(x)\,u$. These formulas will identify a line twist with a
fractional coefficient lattice.

We use $\Delta[x^{-1}]$ for the two-sided localization at the powers
of $x$. Here is the precise reason this localization contains
$\Delta$. With coefficients on the left, the skew relation gives
\[
 hx=x\,\frac{h(x)}x\,h,
 \qquad
 xh=\frac{x}{h(x)}\,h\,x
 \qquad(h\in H).
\]
Both displayed fractions are units of $R$. As in the proof of
Proposition~\ref{geom:skew-filtration}, at coefficient precision
$x^d$ we may refine the group quotient until it detects these units
modulo $x^d$, using the coefficient action modulo $x^{d+1}$.
The two inverse continuous changes of coefficients give
$\Delta x=x\Delta$. Left multiplication by $x$ is injective on
the compatible coefficient systems defining $\Delta$, and the unit
change gives right injectivity as well. Thus $x$ is normal and regular,
so its powers admit this two-sided localization. This coefficient
element is generally noncentral.

We shall compute the graded algebra
\[
 \mathscr E_n=[\widehat T'[n],\widehat T'].
\]
The product convention is composition: $ab$ means first apply $b$
and then apply the appropriately shifted $a$. With our periodicity
convention, an ordinary sheaf map $T'_d\to T'_d(s)$ gives a map
$T'_d[sD]\to T'_d$. The calculation must retain these ordinary maps:
their strict compositions will make matrices of operations into
complexes.

\begin{proposition}[Complete operations and their representatives]
\label{complete:operations}
There is an identification of graded algebras
\begin{equation}
\label{complete:operation-algebra}
 \mathscr E_{sD}=(x^{-s}\Delta)^\Lambda\quad(s\in\Z),
 \qquad
 \mathscr E_n=0\quad(n\notin D\Z).
\end{equation}
Products use left-coefficient skew convolution in $\Delta[x^{-1}]$,
with degrees added. Every element of $\mathscr E_{sD}$ has a unique
compatible system of
ordinary sheaf-map representatives
$T'_d\to T'_d(s)$, and these representatives realize multiplication
strictly. Under \eqref{complete:operation-algebra}, the operation $t$
of degree $D$ is the coefficient $1\in x^{-1}\Delta$; it is central
in $\mathscr E$.

Put
\[
 A_{\mathrm{op}}=(\Delta[x^{-1}])^\Lambda.
\]
With a formal $t$ recording degree $D$, localization identifies
\begin{equation}
\label{complete:localized-operations}
 \mathscr E[t^{-1}]=A_{\mathrm{op}}[t^{\pm1}].
\end{equation}
\end{proposition}

\begin{proof}
We first exclude the extension and inverse-limit terms so that all
operations have ordinary compatible representatives. We then normalize
the line twists and compute composition at every finite precision.

The fiber of $T'\to\widehat T'$ is
$t$-invertible, so Proposition~\ref{complete:fracture} identifies maps
into $\widehat T'$ out of $T'$ or out of $\widehat T'$.
Maps into the limit can consequently be computed from the tower
of targets $T'_d$.

Flat affine adjunction along $f'$ computes
$[T'[n],T'_d]$ on $\mathcal X_k$, with target
$\mathcal O_{\mathcal X_k/(t^d)}$. This target is pushed forward
from the open $v_2\ne0$, so we may work on
$[\Spec k[x]/\Lambda]$. Let $M$ be the restriction of
$(f')^*T'$ to this chart. It is $x$-monic and
\[
 M/x^dM=C_{\mathrm{lc}}(H,R_d).
\]
As an $R_d$-module the latter is the scalar extension of
$C_{\mathrm{lc}}(H,k)$, hence is free. Since $|\Lambda|=q$
is invertible in $k$, averaging a lift over $\Lambda$ proves
equivariant projectivity as well.

Choose an equivariant projective resolution of $M$ of length one
over $k[x]$. Start with a surjection from a free module over the
skew group ring $k[x]\#\Lambda$; its underlying $k[x]$-module is free.
Since $k[x]$ is hereditary, its kernel is projective as a
$k[x]$-module; averaging again promotes this to equivariant projectivity.
Reduction modulo $x^d$ preserves exactness, since $M$ is $x$-monic.
Applying $\Hom$ into $R_du^s$ factors through this reduced
resolution. Its positive cohomology vanishes by the projectivity of
$M/x^dM$ just proved. Periodize the length-one resolution, with the
specified line twists. Its bounded resolution length and projective
terms make it valid for computing maps to twisted periodic complexes:
mapping from it takes acyclics to acyclics. Since $D>2$, the
augmentation degree and the relation degree cannot coincide modulo
$D$. Thus the only possible groups are the ordinary maps in degrees
$sD$ and their first extensions; the latter vanish. In particular,
\begin{equation}
\label{complete:finite-level-maps}
 [T'[sD],T'_d]
   =\Hom\bigl(T',T'_d(s)\bigr),
 \qquad
 [T'[n],T'_d]=0\quad(n\notin D\Z).
\end{equation}

These ordinary Hom groups have surjective transition maps. Indeed,
a map at level $d$ can be viewed as a functional on
$M/x^{d+1}M$ with values in $R_du^s$; equivariant projectivity of
$M/x^{d+1}M$ lifts it through
$R_{d+1}u^s\twoheadrightarrow R_du^s$. The Milnor exact sequence
therefore has zero derived-limit term. Each ordinary map in
\eqref{complete:finite-level-maps} kills $t^d$ and factors uniquely
through $T'_d$. Compatibility in $d$ gives actual maps of towers of
pure complexes, and composition is their ordinary composition. This
factorization is a statement about sheaf maps; it does not assert
full faithfulness of closed pushforward on derived categories.

We next identify these representatives and their product. Evaluation
at the identity coset in coinduction identifies a map with a
$\Lambda$-equivariant $R_d$-linear functional
\[
 C_{\mathrm{lc}}(H,R_d)\longrightarrow R_du^s.
\]
Without the twist, dualizing the functions on each finite group
quotient gives an $R_d$-valued measure, with masses added under
passage to a smaller quotient. Taking both group and coefficient
limits gives $\Delta$. To include the twist, use
\begin{equation}
\label{complete:twist-normalization}
 R_du^s\ \cong\ x^{-s}R/x^{d-s}R,
 \qquad u^s\longmapsto x^{-s}.
\end{equation}
This is compatible with the group action: the formula
$g(u)=(x/g(x))u$ of Proposition~\ref{geom:thickened-charts}
becomes the ordinary action on the fractional coefficient
$x^{-s}$. For the tame action, $\lambda(u)=\lambda u$ follows from
the invariance of $t=xu$, and agrees with
$\lambda(x^{-s})=\lambda^s x^{-s}$. Thus the normalization is
equivariant for every positive or negative $s$. Equivariance of the
functional is therefore invariance of its measure under the coefficient
and conjugation action of $\Lambda$. This gives the groups in
\eqref{complete:operation-algebra}.

For clarity, write integrals in left-coefficient notation. They are
finite sums after passage to a sufficiently deep coefficient and
group quotient, and their values are compatible as those quotients
are refined. If $b$ is the measure describing a map and $F$ its
input function, equivariance reconstructs its value at $h\in H$ as
\begin{equation}
\label{complete:coinduced-operator}
 h\left(\int_z b(z)(h^{-1}F)(z)\right)
       =\int_z h(b(z))F(hz).
\end{equation}
The same formula applies to fractional coefficient lattices using
\eqref{complete:twist-normalization}. In our diagonal semilinear
coinduction coordinates the projection formula is the pointwise map
\[
 C_{\mathrm{lc}}(H,R_d)\otimes L_j
 \longrightarrow C_{\mathrm{lc}}(H,L_j),\qquad
 F\otimes l\longmapsto\bigl(h\longmapsto F(h)l\bigr),
 \qquad L_j=R_du^j.
\]
It is equivariant: applying $g$ to either side gives
$g(F(g^{-1}h))g(l)$ at $h$. In particular the second factor here
is $l$, not $h(l)$. Tensoring both sides of a map by the $j$th line
twist therefore multiplies both its input
function and its output at the identity by $x^{-j}$. Its normalized
measure is unchanged. This is the $R$-linear extension of the
evaluation functional to the shifted lattice; it is not a claim
that $x$ commutes with the resulting skew operator on sections.

Apply the map $b$ first and then the map $a$. Evaluation of their
composite, using \eqref{complete:coinduced-operator}, is
\[
 \int_h\int_z a(h)h(b(z))F(hz),
\]
which is precisely left-coefficient skew convolution $ab$. All these
identities hold at each finite level, including for negative twists.
For example, if the degrees of $a$ and $b$ are $sD$ and $jD$,
changing their representatives respectively by coefficients in
$x^{d-s}R$ and $x^{d-j}R$ changes the product only by
$x^{d-s-j}R$. This follows because every $h$ preserves $x$-adic
order. Group quotients may always be refined to detect the relevant
units $h(x)/x$ and their inverses to the required coefficient order.
Thus composition and reduction are strictly compatible, as asserted.

On invariant sections the input function is the orbit function of
its value $f$ at the identity. The action consequently becomes
\begin{equation}
\label{complete:action-on-sections}
 a(f)=\int_H a(h)h(f).
\end{equation}
Multiplication by $t=xu$ normalizes to the coefficient $1$, proving
the assertion about $t$. In particular it is central in the graded
operation algebra, although $x$ need not be central in $\Delta$.
Finally, inverting this degree-$D$ inclusion allows arbitrarily large
fractional coefficient lattices in every fixed degree. Their union
is $\Delta[x^{-1}]$, and taking $\Lambda$-invariants commutes with
this union. This proves \eqref{complete:localized-operations}.
\end{proof}

The same finite-level calculation gives the vanishing of self-maps to
the $2$-shift for all three characteristic-$p$ chart sheaves. The
assertion includes every shift of each sheaf.

\begin{corollary}
\label{complete:degree-two}
For each $j\in\Z$ there are no self-maps to the $2$-shift of
$\widehat T'[j]$, $\widehat T_k[j]$, or $\widehat T_0[j]$.
\end{corollary}

\begin{proof}
For $\widehat T'$ this is immediate from
\eqref{complete:operation-algebra}. The finite-level calculation
\eqref{complete:finite-level-maps} applies equally to the full-center
chart sheaf $T_k$, replacing the functions on $H$ by the functions
for $G_{\mathrm{up}}$. They are again free over $R_d$ and
equivariantly projective for the tame subgroup $\Lambda$. Hence
the same extension and Milnor terms vanish.

For $T_0$, use extension of scalars
$e^*:\mathbf D(\mathcal M_0)\to\mathbf D(\mathcal M_k)$.
Since $k/\F_p$ is finite separable, trace duality identifies this
extension functor also as a right adjoint to restriction. It
preserves homotopy limits, and therefore commutes with completion.
Moreover the unit $F\to e_*e^*F$ is naturally split as an
$\F_p$-linear map, so extension of scalars detects zero maps.
The asserted vanishing follows from that for $\widehat T_k$.
\end{proof}

\subsection{Finite constructions from the complete generator}

The operations just computed act strictly on the entire quotient
tower. We now turn this action into a description of all finite
constructions from $\widehat T'$. The important point is that the
relations in a differential matrix will hold as equalities of maps
at every level, so its cones and sections can be calculated from the
same operators.

Give $\mathscr E$ its homological grading and zero differential. A
\emph{finite-cell right dg $\mathscr E$-module} is a finite iterated
extension of shifts of the free right module $\mathscr E$. It can be
represented by a finite graded-free module with a filtration whose
successive quotients are shifted free modules. A perfect right module
is a retract of a finite-cell module in the derived category; their
category is $\Perf(\mathscr E)$.

The description of an algebraic thick subcategory by perfect modules
over a dg endomorphism algebra is a general form of dg Morita theory;
compare Keller \cite[Section~3.6 and Theorem~3.8(b)]{Keller2006}.
Here we use the computed graded algebra $\mathscr E$ itself with zero
differential. Its strict action on the finite reductions, proved in
Proposition~\ref{complete:operations}, is what permits this choice.

\begin{proposition}[Strict realization]
\label{complete:realization}
There is an exact fully faithful functor
\[
 \mathcal R:\Perf(\mathscr E)\longrightarrow\mathbf D(\mathcal N)
\]
from perfect right dg modules which sends $\mathscr E$ to
$\widehat T'$ and identifies its image with
$\Thick(\widehat T')$. For a finite-cell module $\mathsf P$, its
image is the homotopy limit of an actual compatible tower $Y_d$ of
complexes whose terms are finite sums of shifts of $T'_d$.
\end{proposition}

\begin{proof}
First let $\mathsf P$ be a finite-cell module. Choose a homogeneous
basis of $\mathsf P$ compatible with its cell
filtration. Replace a basis vector $e_b$ of degree $b$ by the shifted
pure complex $T'_d[b]$ at every level $d$. A degree-$j$ right-module
matrix entry
\[
 e_b\longmapsto e_{b'}a,
 \qquad a\in\mathscr E_{b+j-b'},
\]
is represented in chain degree $m$ by
\begin{equation}
\label{complete:matrix-degrees}
 a:(T'_d)_{m-b}\longrightarrow(T'_d)_{m+j-b'}.
\end{equation}
Use the unique representatives of
Proposition~\ref{complete:operations} for all entries, including the
differential entries, where $j=-1$. The source and target use the same
line-twist periodicity identifications in every degree and at every
level $d$. On the normalized charts these are the powers of the descent
unit for $u$ already used in \eqref{complete:twist-normalization}.

Strict composition makes the matrix relation $d^2=0$ an equality of
actual maps. All nonzero degrees of $\mathscr E$ are even, so no
coefficient sign enters this relation. For a degree-$j$ matrix $f$,
the Hom differential is
\[
 d_{\mathsf Q}f-(-1)^j f d_{\mathsf P}.
\]
Formula~\eqref{complete:matrix-degrees} respects this differential,
and therefore respects chain maps, chain homotopies, the signed shift,
and ordinary mapping cones. It also makes all these complexes and maps
compatible under the transitions $d+1\to d$. Every matrix entry uses
its compatible tower map, so its representative is fixed at all levels
at once.

The order of multiplication agrees with the right-module convention.
An endomorphism of the free right module determined by $1\mapsto a$
sends $z\mapsto az$, by left multiplication. Composing the
endomorphism for $b$ and then the one for $a$ gives $ab$, the same
order as the operators in Proposition~\ref{complete:operations}.
Tensoring a right module with a left module consequently uses the
left action of these same operators.

Take the homotopy limit of the constructed tower. Limits are additive
and exact in the stable category, so the realization of a levelwise
cone is the cone of the realized map. Compatible chain homotopies
induce equal maps. A change of homogeneous basis gives isomorphic
towers through its change-of-basis matrices. We obtain an exact functor
on the homotopy category of finite cells. Mapping a shifted free module
to an acyclic module gives an acyclic Hom complex, and induction along
a finite free-cell filtration preserves this property. Finite cells are
therefore homotopically projective, so their homotopy maps compute their
derived maps.

On the free generator and every shift, the functor induces precisely
the isomorphisms of Proposition~\ref{complete:operations}. Fix a free
source and use the long exact sequences of maps along a cell filtration
of the target to obtain the isomorphism for every finite-cell target.
Then apply the long exact sequences along a cell filtration of the
source. This proves full faithfulness for all finite cells.
The target derived category is idempotent complete, so the functor
extends fully faithfully to retracts. Its image is closed under shifts,
finite sums, cones, and retracts and contains $\widehat T'$. Conversely,
every object in its image is constructed by those operations from
$\widehat T'$. Its image is therefore exactly
$\Thick(\widehat T')$.
\end{proof}

For example, if $a\in\mathscr E_0$, the two-cell module
$\cofib(a:\mathscr E\to\mathscr E)$ realizes
$\cofib(a:\widehat T'\to\widehat T')$. The module map is left
multiplication on the free right module. We next compute the sections
of this cone using the same left operator on a coefficient lattice;
its kernel and cokernel are thus sections data of an actual completed
object.

\subsection{Sections and finite reductions}

Let $\mathsf L$ be the graded left $\mathscr E$-module
\begin{equation}
\label{complete:section-module}
 \mathsf L_{sD}=(x^{-s}R)^\Lambda\quad(s\in\Z),
 \qquad \mathsf L_n=0\quad(n\notin D\Z),
\end{equation}
with zero differential and action
\eqref{complete:action-on-sections}. Its finite reductions are the
graded left modules
\[
 (\mathsf L_d)_{sD}=(x^{-s}R/x^{d-s}R)^\Lambda\quad(s\in\Z),
 \qquad (\mathsf L_d)_n=0\quad(n\notin D\Z),
\]
and $\mathsf L=\varprojlim_d\mathsf L_d$ degree by degree. Set
\[
 \Phi=(k((x)))^\Lambda=k((x^q)).
\]
Here $t$ acts on normalized sections by the inclusion of one
fractional lattice in the next. In particular
$\mathsf L[t^{-1}]=\Phi[t^{\pm1}]$, with $t$ recording degree
$D$. For a complete object $Y$ write $W(Y)=\pi_*LY$.

\begin{proposition}[Section formulas]
\label{complete:sections}
If $Y=\mathcal R(\mathsf P)$ for a perfect right dg
$\mathscr E$-module $\mathsf P$, then naturally
\begin{equation}
\label{complete:tensor-formulas}
 \begin{split}
 \pi_*Y&=H_*\bigl(\mathsf P\otimes^{\mathbf L}_{\mathscr E}
                           \mathsf L\bigr),\\
 W(Y)&=H_*\bigl(\mathsf P\otimes^{\mathbf L}_{\mathscr E}
                           \Phi[t^{\pm1}]\bigr).
 \end{split}
\end{equation}
For a finite-cell model the ordinary cellular tensor computes each
derived tensor. The actions in these formulas are the left actions
of \eqref{complete:action-on-sections}.
\end{proposition}

\begin{proof}
Suppose first that $\mathsf P$ is a finite-cell module and write
$Y_d$ for its realization tower. Each term of $Y_d$ is a finite
sum of twists of $T'_d$. Those twists are sections-acyclic: flat
affine adjunction reduces sections to the nilpotent chart of
$\mathcal X_k$, where taking $\Lambda$-invariants is exact.
The bounded derived-sections calculation from
Proposition~\ref{geom:cohomological-bounds} therefore computes
$\pi_*Y_d$ by the termwise sections complex. By
\eqref{complete:action-on-sections}, that complex is
$\mathsf P\otimes_{\mathscr E}\mathsf L_d$. This identifies
differentials and maps, as well as the underlying groups.

Write $C_d=\mathsf P\otimes_{\mathscr E}\mathsf L_d$. In each fixed
degree its terms, cycles, boundaries, and homology are finite groups:
$k$ is finite, the fractional lattices modulo $x^d$ are finite
dimensional, and there are only finitely many cells. Each of these
towers therefore satisfies the Mittag--Leffler condition, so its first
derived inverse limit vanishes. Taking inverse limits in the exact
sequences defining cycles, boundaries, and homology shows that
$H_*(\varprojlim_d C_d)=\varprojlim_d H_*(C_d)$. The Milnor sequence
for $Y=\holim_dY_d$ and the finite graded-free underlying module of
$\mathsf P$ now give
\[
 \pi_*Y
 =H_*\left(\varprojlim_d
       \mathsf P\otimes_{\mathscr E}\mathsf L_d\right)
 =H_*\bigl(\mathsf P\otimes_{\mathscr E}\mathsf L\bigr).
\]
The tensor-limit interchange is degreewise and involves only finitely
many summands. Finite-cell right modules are homotopically flat: after
tensoring their cell filtration with an acyclic left module, every
successive quotient is a shift of that acyclic module, so induction
makes the full tensor acyclic. Thus the ordinary cellular tensor
computes the derived tensor in the first formula.

Now take the $t$-colimit. Proposition~\ref{complete:fracture} identifies
it with the localized sections. Filtered colimits are exact, and
$\mathsf L[t^{-1}]=\Phi[t^{\pm1}]$, which gives the second formula
with the same homotopical flatness justification. The identifications
are natural for finite-cell maps. Applying them to idempotent summands
proves both formulas for all perfect modules.
\end{proof}

The topology of the overlap vector spaces will use these complete
section groups. To make the topology on $\pi_iY$ intrinsic, we identify
the tower just constructed with the reductions of the complete object.

\begin{lemma}[The tower is the reduction tower]
\label{complete:reductions}
For a finite-cell realization $Y=\holim_eY_e$ there are natural
equivalences
\[
 Y/t^d\simeq Y_d\qquad(d\geq1).
\]
Under these equivalences, the tower projection is the intrinsic
reduction map. The profinite topology on $\pi_iY$ obtained from the
finite-level sections calculation is consequently its inverse-limit
topology from $\pi_i(Y/t^d)$.

If $Y'$ is a retract of $Y$, let $e_d$ be the idempotent of $Y_d$ in
the derived category induced by the idempotent defining this retract.
Then $Y'/t^d$ is the summand of $Y_d$ split by $e_d$, and
\[
 \pi_iY'=\varprojlim_d\operatorname{im}(\pi_i e_d)
\]
with its intrinsic profinite topology.
\end{lemma}

\begin{proof}
Multiplication by $t^d$ is strictly zero on $Y_d$, because it is
strictly zero on every shifted $T'_d$ in its terms. For $e\geq d$,
map the ordinary cone of $t^d:Y_e[dD]\to Y_e$ to $Y_d$ by the
tower transition on the unshifted summand and zero on the shifted
summand. Strict annihilation makes this a chain map. These maps
are compatible in $e$ and natural for the matrix maps defining our
finite-cell constructions. Exactness of homotopy limits gives a
map $Y/t^d\to Y_d$.

On the free generator this is the equivalence
$\widehat T'/t^d\simeq T'/t^d=T'_d$ of
Proposition~\ref{complete:fracture}, represented by the ordinary
quotient. The comparison respects mapping cones and shifts, with the
usual sign isomorphism for a shifted cone. Induction on the finite cell
filtration proves the equivalence for $Y$ and identifies the tower
projection with its intrinsic reduction map.

The groups $\pi_iY_d$ are finite by the section calculation, and their
first derived inverse limits vanish. Hence
$\pi_iY=\varprojlim_d\pi_iY_d$ is profinite, and the natural
equivalences just proved identify this topology with the intrinsic
one. For a retract $Y'$ of $Y$, naturality sends its defining derived
idempotent to an idempotent $e_d$ on every $Y_d$. Reduction preserves
the retract, so $Y'/t^d$ is the summand split by $e_d$; this statement
takes place in the derived category and requires no strict idempotent
representative. On the finite group $\pi_iY_d$, the map $\pi_i e_d$
is an ordinary idempotent. Taking its image commutes with the inverse
limit, because the associated image and kernel decompositions are
compatible with the transition maps. This gives the displayed formula
for $\pi_iY'$ and its intrinsic reduction topology.
\end{proof}

\subsection{Compact and discrete parts of the overlap}

We have described the complete objects built from $\widehat T'$ by
finite algebraic constructions. We now study the topology of the overlap
vector spaces. When an object has finite sections homotopy in every degree
before completion, its open image occupies a discrete cocompact subgroup
of the overlap. For completions of objects with this finiteness, the
topology is intrinsic, and every map between them induces continuous
maps on the overlap vector spaces.

\needspace{26\baselineskip}
\begin{proposition}[Intrinsic overlap topology]
\label{complete:topology}
Every $P\in\Thick(T')$ has finite $\pi_iP$ for every integer $i$.
Now let $P\in\mathbf D(\mathcal N)$ have finite $\pi_iP$ for every
integer $i$. For each integer $i$, the group $\pi_i\widehat P$ is
naturally profinite, with the intrinsic inverse-limit topology from the
finite groups $\pi_i(\widehat P/t^d)$. The vector space
$W_i(\widehat P)=\pi_iL\widehat P$ has a locally compact Hausdorff
$k$-vector-space topology obtained as the increasing union of the compact
open images of $\pi_{i+sD}\widehat P$, for $s\geq0$. The image of
$\pi_iLP$ is discrete and cocompact in this space.

If $Q\in\mathbf D(\mathcal N)$ also has finite $\pi_jQ$ for every
integer $j$, then every map $f:\widehat P\to\widehat Q$ induces
continuous maps $\pi_i\widehat P\to\pi_i\widehat Q$ and
$W_i(\widehat P)\to W_i(\widehat Q)$ for every integer $i$, in the
topologies just described. No global map $P\to Q$ inducing $f$ is
required. If $\widehat P\simeq\mathcal R(\mathsf P)$ for a perfect
right dg $\mathscr E$-module $\mathsf P$, then under this identification
the natural identifications of Proposition~\ref{complete:sections} and
Lemma~\ref{complete:reductions} match the profinite topology on
$\pi_i\widehat P$ with the one obtained from the realization's finite-level
sections and intrinsic reductions. Consequently they match the compact
images and the increasing-union topology on $W_i(\widehat P)$.
\end{proposition}

\begin{proof}
For $T'$, flat affine adjunction identifies sections-homotopy with
the cohomology of twists on
$\mathcal X_k=\mathcal P_k(1,q)$. The two-variable principal-open
\v Cech complex computes it, also in the twisted periodic category.
Weight by weight, its $H^0$ is spanned by nonnegative monomials of
that weight and its $H^1$ by Laurent monomials with both powers
negative. Both sets are finite in a fixed weight because the two
weights $1$ and $q$ are positive. Thus every $\pi_iT'$ is finite.
Degreewise finiteness is preserved by shifts, finite sums, triangles,
and retracts, proving the first assertion.

Now let $P$ satisfy the stated degreewise finiteness hypothesis.
The cofiber sequence for $t^d$ makes every $\pi_i(P/t^d)$ finite.
The Milnor sequence therefore gives
\begin{equation}
\label{complete:compact-homotopy}
 \pi_i\widehat P=\varprojlim_d\pi_i(P/t^d),
\end{equation}
with zero first derived-limit term. Give this group its compact
Hausdorff inverse-limit topology. The equivalence
$P/t^d\simeq\widehat P/t^d$ is compatible with the projections:
both maps $\widehat P\to P/t^d$ agree after precomposition by
$P\to\widehat P$, and this determines a map uniquely because
its fiber is $t$-invertible whereas the target is complete. This
proves that \eqref{complete:compact-homotopy} describes the intrinsic
reduction topology.

The fracture square gives a long exact sequence with the map
\begin{equation}
\label{complete:joint-overlap-map}
 \pi_i\widehat P\oplus\pi_iLP
        \longrightarrow W_i(\widehat P).
\end{equation}
Its kernel is a quotient of $\pi_iP$ and its cokernel embeds in
$\pi_{i-1}P$. In particular both are finite. The same holds in
every shifted degree. Write $C_s=\pi_{i+sD}\widehat P$ and let
$K_s$ be its image in $W_i(\widehat P)$ under the localization
identifications. The kernel $C_s\to K_s$ is finite, by applying
\eqref{complete:joint-overlap-map} in degree $i+sD$. Give $K_s$
its compact Hausdorff quotient topology.

The maps $C_s\to C_{s+1}$ are continuous by naturality of the
intrinsic reduction maps. Their cokernels are finite, by the
multiplication-cofiber sequence and
$\widehat P/t\simeq P/t$. Thus the resulting inclusions
$K_s\subset K_{s+1}$ are continuous and of finite index. A
continuous injection from a compact space to a Hausdorff space is
a closed embedding; here finite index also makes it open. By
\eqref{complete:localization-sections}, the union of these compact
open subgroups is exactly $W_i(\widehat P)$. Give the union its
corresponding topology. This is a locally compact Hausdorff
topological $k$-vector space.

Let $B$ be the image of $\pi_iLP$. Its intersection with $K_0$ is
finite, because a common image determines a pair in the finite
kernel of \eqref{complete:joint-overlap-map}. In the Hausdorff
group $K_0$, choose a neighborhood of zero avoiding the nonzero
points of this finite intersection. It meets $B$ only in zero,
so $B$ is discrete. A discrete subgroup of a Hausdorff topological
group is closed. The quotient by $K_0+B$ is finite, by the
finite cokernel of \eqref{complete:joint-overlap-map}; hence
$W_i(\widehat P)/B$ is covered by finitely many translates of the
compact image of $K_0$. Thus $B$ is cocompact.

Finally, let $Q$ and $f:\widehat P\to\widehat Q$ be as in the
statement. The map $f$ induces maps of the intrinsic finite reduction
groups, so it is continuous on the profinite homotopy groups. After
localization its restriction to each compact stage is continuous, so it
is continuous on the unions just constructed. No global map $P\to Q$
inducing $f$ is needed. The finite-level section calculation of
Proposition~\ref{complete:sections} and the natural reduction equivalences
of Lemma~\ref{complete:reductions} identify the profinite topologies for
the perfect realization of $\widehat P$ in the statement. These
identifications commute with $t$, so they identify each compact image and
the union topology after localization. This includes the case where
$\mathsf P$ is a retract of a finite-cell module.
\end{proof}
